\documentclass[openany,12pt,UTF8]{ctexbook}
\usepackage[a4paper,twoside,margin=2.5cm]{geometry}
% ============================
% §1 数学包
% ============================
\usepackage{amsmath}                                    % 排版数学公式
\usepackage{latexsym}
\usepackage{amsfonts}                                   % AMS 数学字体（Computer Modern 补充符号字体）
\usepackage{amssymb}                                    % AMS 数学符号命令
\usepackage{amscd}                                      % 交换图环境（如不用可删除）
\usepackage{mathrsfs}                                   % 数学 RSFS 书写字体
\usepackage{bm}                                         % 数学黑体（\bm 命令）
\usepackage{physics}                                    % 物理公式排版
\usepackage[math-style=ISO]{unicode-math}               % 使用 ISO 数学风格

% 计数器依赖 section
\numberwithin{equation}{section}
\numberwithin{table}{section}
\numberwithin{figure}{section}

% ============================
% §2 排版格式包
% ============================
\usepackage{graphicx}                                   % 支持插图
\usepackage{color,xcolor}                               % 支持彩色
\usepackage{titlesec}                                   % 设置章节格式
\usepackage{tcolorbox}                                  % 彩色文本框
\usepackage{enumerate}                                  % 更改 enumerate 环境格式
\usepackage[perpage]{footmisc}                          % 脚注每页重新编号

% 设置 paragraph 为只有一级编号（从1开始）
\renewcommand{\theparagraph}{\arabic{paragraph}}
\makeatletter
\@addtoreset{paragraph}{section}                        % 每节开始时重置 paragraph 计数器
\makeatother

% 设置 paragraph 为 block 风格，自动换行
\titleformat{\paragraph}[block]
  {\normalfont\normalsize\bfseries}{\theparagraph.}{0.5em}{}

% 设置标题后换行（after-sep 设为 \newline）
\titlespacing*{\paragraph}{0pt}{3.25ex plus 1ex minus .2ex}{0pt}

% ============================
% §3 代码与算法包
% ============================
\usepackage{minted}                                     % 代码高亮
\usepackage[linesnumbered,ruled,vlined]{algorithm2e}    % 算法排版

% ============================
% §4 表格与图表包
% ============================
\usepackage{diagbox}                                    % 表格斜线表头
\usepackage{tabularx}                                   % 自动调整列宽
\usepackage{multirow}                                   % 合并行
\usepackage{subcaption}                                 % 子图和子标题
\usepackage{minipage-marginpar}                         % minipage 中使用 marginpar
\usepackage{float}                                      % 提供 float 浮动环境
\usepackage{booktabs}                                   % 三线表（\toprule、\midrule、\bottomrule）
\usepackage{listings}                                   % 代码排版（非高亮）

% ============================
% §5 引用与超链接包
% ============================
\usepackage{hyperref}

% 改变超链接颜色
\hypersetup{
    colorlinks=true,
    linkcolor=blue,
    filecolor=blue,
    urlcolor=blue,
    citecolor=cyan,
}

% autoref 中文引用名
\def\equationautorefname{式}
\def\footnoteautorefname{脚注}
\def\itemautorefname{项}
\def\figureautorefname{图}
\def\tableautorefname{表}
\def\partautorefname{篇}
\def\appendixautorefname{附录}
\def\chapterautorefname{章}
\def\sectionautorefname{节}
\def\subsectionautorefname{小小节}
\def\subsubsectionautorefname{小小小节}
\def\paragraphautorefname{段落}
\def\subparagraphautorefname{子段落}
\def\FancyVerbLineautorefname{行}
\def\theoremautorefname{定理}

% ============================
% §6 标题层级与目录设置
% ============================
\setcounter{tocdepth}{3}
\setcounter{secnumdepth}{4}

% ============================
% §7 自定义命令
% ============================
% 角标引用文献和交叉引用
\newcommand{\scite}[1]{\textsuperscript{\cite{#1}}}
\newcommand{\sref}[1]{\textsuperscript{\ref{#1}}}
% boldsymbol 简写别名（等价于 bm 包的 \bm）
\newcommand{\bs}[1]{\boldsymbol{#1}}
% 罗马数字
\newcommand{\romannumber}[1]{\uppercase\expandafter{\romannumeral#1}}
% 红色警示文本
\newcommand{\alert}[1]{\textcolor{red}{#1}}

% 彩色文本框，宽度自适应版心
\newcommand{\tbox}[1]{
  \begin{center}
    \begin{tcolorbox}[colback=gray!10,
        colframe=black,
        width=\linewidth,
        arc=1mm, auto outer arc,
        boxrule=0.5pt,
      ]
      {#1}
    \end{tcolorbox}
  \end{center}
}

% 法语强调下划线
\newcommand{\fri}[1]{\underline{#1}}

\begin{document}
\frontmatter
\begin{titlepage}
	\centering
	\scshape
	\vspace*{1.5\baselineskip}

	\rule{13cm}{1.6pt}\vspace*{-\baselineskip}\vspace*{2pt}
	\rule{13cm}{0.4pt}

		\vspace{0.75\baselineskip}
	{
        \Huge
			《红枫 2026 无人集群自主协同\\[0.3em]智能算法挑战赛 技术报告》 \\
    }
		\vspace{0.75\baselineskip}
	\rule{13cm}{0.4pt}\vspace*{-\baselineskip}\vspace{3.2pt}
	\rule{13cm}{1.6pt}

		\vspace{1.75\baselineskip}
	{\large
		参赛队伍：[队伍名称] \\
		\vspace*{1.2\baselineskip}
		参赛成员：[成员姓名] \\
		\vspace*{1.2\baselineskip}
		指导教师：[教师姓名]（如有） \\
		\vspace*{1.2\baselineskip}
		单位：[学校/单位] \\
	}
	\vfill

	{\large \today}

\end{titlepage}

\tableofcontents
\mainmatter

%%% ============================================================ %%%
%%% 正文开始 %%%
%%% ============================================================ %%%

\chapter{研究背景与竞赛任务}

\section{无人集群 ISR 任务概述}

% 当前无人集群在情报、监视、侦察领域的应用背景与核心挑战

\section{仿真平台与竞赛规则}

% 平台架构（C++ 底层 + Python SDK）、信息隔离机制、评分体系

\section{三道赛题的递进关系}

% 赛题一（单机搜索跟踪）→ 赛题二（多机协同识别）→ 赛题三（对抗集群全域搜索）

\chapter{方法论与算法设计}

\section{总体架构设计}

% 分层架构：Agent 层 → 状态机层 → 算法模块层 → 计算服务层（Rust）

\section{搜索算法}

% 区域覆盖搜索策略：螺旋搜索、网格分配、概率图引导搜索

\section{目标跟踪与状态估计}

% 云台 LOS 瞄准、EMA/卡尔曼滤波、bearing-only 跟踪的可观测性分析

\section{诱饵鉴别算法}

% 运动学一致性检验：速度估计、位置散布分析、多帧判据

\section{多机协同与任务分配}

% 通信协议设计、K=2/K=3 分布式共识、Greedy/拍卖分配算法

\section{威胁感知与规避}

% SAM 防空区规避（高度策略+水平绕行）、动态干扰区感知与广播预警

\chapter{核心算法实现}

\section{Rust 计算服务}

% PyO3 集成、卡尔曼滤波器、路径规划（A*）、几何计算加速

\section{Python 算法模块}

% 按功能划分的模块结构：search/tracking/discrimination/coordination/threat/estimation

\section{通信协议设计}

% 消息格式（50 字节约束）、4 种消息类型（T/D/A/J）、编码效率

\section{状态机设计}

% ACQUIRE → VERIFY → TRACK → SEARCH 状态转换逻辑

\chapter{实验结果与分析}

\section{赛题一实验}

% 搜索效率、跟踪精度、目指上报 RMSE

\section{赛题二实验}

% 诱饵鉴别准确率、K=2 协同效率、全歼时间

\section{赛题三实验}

% 存活率、摧毁率、威胁规避效果、K=3 协同分配效率

\section{消融实验}

% 各算法模块的贡献分析

\chapter{创新点与应用价值}

\section{技术创新点}

% 列举 3-5 个具体创新点

\section{应用前景}

% 应急搜救、边境巡查、侦察等实际场景应用价值

\chapter{总结与展望}

\section{工作总结}

\section{不足与改进方向}

%%% 正文结束 %%%
%%% ============================================================ %%%

\end{document}
